% Ad Familiares 6.22 (ad-familiares-06-22) --- English translation
% Translated by Claude (Anthropic) from the Latin (Perseus).
% Style guide: STYLE.md.

\ciceroLetterOpener{CICERO DOMITIO.}

\ciceroSection{1} It was not the fact that you had sent me no letter that deterred me from writing to you, after you had come into Italy; rather, in such great evils I could find nothing to promise --- I myself being in want of everything --- nothing to advise --- when counsel was lacking for my own case --- nothing in the way of consolation to bring. Though these matters are no better now, and indeed by far more desperate, still I have preferred my letter to be empty rather than absent altogether.

\ciceroSection{2} If I understood that you had taken upon yourself, for the sake of the state, more of a charge than you could have made good, I would still, by whatever means I could, urge you to that mode of living which was on offer and which lay open. But since you have set, for a course chosen well and bravely, the limit which fortune herself would have wished to be the boundary of our struggle, I beg and beseech you, by our long-standing tie and bond, and by my very great goodwill toward you and yours toward me which equals it, that for our sake, for your mother's sake, for your wife's, for all those to whom you are and always have been most dear, you keep yourself safe; that you take thought for your own preservation and that of yours who depend upon you; that the lessons you have learned, the things which from your youth onward you have most beautifully grasped in memory and understanding from the wisest of men, you put to use at this time; and that the loss of those whom, joined to you by the highest goodwill and by countless services, you have now lost, you bear --- if not with an even mind, then at least with a brave one.

\ciceroSection{3} What I can do, I do not know, or rather I feel I can do little; this much, however, I promise you: whatever I judge to bear upon your safety and your standing, I will set my hand to with as much zeal as you have always shown, in zeal and in service, in my own affairs. This intention of mine I have conveyed to your mother, that excellent woman who loves you so deeply. If you write me anything, I shall act as I understand you to wish; but if you write less, I shall still, with the greatest zeal and care, attend to everything I judge to be of use to you. Farewell.
